\documentclass{beamer}
\usetheme{Madrid}
\usecolortheme{default}
\usepackage{tikz}

\title{Analog-1.1.60\\Control Systems\\Root Locus and Stability}
\author{NAMASWI V-EE25BTECH11060}
\date{\today}

\begin{document}

\begin{frame}
    \titlepage
\end{frame}

%------------------------------------------------

\begin{frame}{Question}

A transfer function with unity DC gain has three poles at $-1,-2,-3$ and no finite zeros. A plant is connected with a proportional controller of gain $K$ in a unity feedback configuration.

\vspace{0.3cm}

(a) The Transfer Function is 

\begin{columns}[t]
    \begin{column}{0.45\textwidth}
        \begin{enumerate}[i)]
            \item $\frac{s}{(s-1)(s-2)(s-3)}$
            \item $\frac{6}{(s+1)(s+2)(s+3)}$
        \end{enumerate}
    \end{column}
    
    \begin{column}{0.45\textwidth}
        \begin{enumerate}[i)]
            \setcounter{enumi}{2}
            \item $\frac{s}{(s+1)(s+2)(s+3)}$
            \item $\frac{6}{(s-1)(s-2)(s-3)}$
        \end{enumerate}
    \end{column}
\end{columns}

\vspace{0.2cm}

(b) If root locus passes through $\pm j\sqrt{11}$, find maximum $K$

\begin{columns}[t]
    \begin{column}{0.45\textwidth}
        \begin{enumerate}[i)]
            \item $\sqrt{11}$
            \item $6$
        \end{enumerate}
    \end{column}
    
    \begin{column}{0.45\textwidth}
        \begin{enumerate}[i)]
            \setcounter{enumi}{2}
            \item $10$
            \item $6\sqrt{11}$
        \end{enumerate}
    \end{column}
\end{columns}

\end{frame}

%------------------------------------------------


\begin{frame}{Understanding the Terminology}

\textbf{Plant/System:} A system which takes an input signal and produces an output signal.

\vspace{0.2cm}

\textbf{Transfer Function:}
\begin{equation}
G(s) = \frac{\text{Output}}{\text{Input}}
\end{equation}

% It tells how input is transformed into output in the Laplace domain.

\vspace{0.2cm}

\textbf{Poles:} Values of $s$ that make the denominator zero.

\textbf{Zeros:} Values of $s$ that make the numerator zero.

\vspace{0.2cm}

\textbf{Gain $K$:} A constant that controls how strong the system response is.

\vspace{0.2cm}

\textbf{Unity Feedback:} Output is directly fed back (no change), i.e., feedback = 1.

\end{frame}



%------------------------------------------------

 %------------------------------------------------

\begin{frame}{How the System Works}

\begin{center}
\begin{tikzpicture}[auto, node distance=2cm, >=latex']

\node [coordinate] (input) {};
\node [circle, draw, minimum size=1cm] (sum) at (2,0) {$\sum$};
\node [rectangle, draw, minimum width=2cm, minimum height=1cm] (block) at (5,0) {$KG(s)$};
\node [coordinate] (output) at (8,0) {};
\node [coordinate] (feedback) at (5,-2) {};

\draw[->] (input) -- node[above] {$R(s)$} (sum);
\draw[->] (sum) -- node[above] {$E(s)$} (block);
\draw[->] (block) -- node[above] {$Y(s)$} (output);
\draw[->] (output) |- (feedback);
\draw[->] (feedback) -- node[left] {$-$} (sum);

\end{tikzpicture}
\end{center}

\end{frame}

%------------------------------------------------

\begin{frame}{Step 1: Error Signal}

At the summing junction:

\begin{equation}
E(s) = R(s) - Y(s)
\end{equation}

\vspace{0.3cm}

\textbf{Meaning:}
\begin{itemize}
    \item $R(s)$ = Input(what we want)
    \item $Y(s)$ = Output(what we get)
    \item $E(s)$ = difference (error)
\end{itemize}

\end{frame}

%------------------------------------------------

\begin{frame}{Step 2: Output of the System}

The block $KG(s)$ takes $E(s)$ as input.

\vspace{0.3cm}

Using:

\begin{equation}
\text{Output} = \text{Transfer Function} \times \text{Input}
\end{equation}

\begin{equation}
Y(s) = KG(s)\,E(s)
\end{equation}

\vspace{0.3cm}

Substitute $E(s)$:

\begin{equation}
Y(s) = KG(s)\big(R(s) - Y(s)\big)
\end{equation}

\end{frame}

%------------------------------------------------

\begin{frame}{Step 3: Final Transfer Function}

Rearranging:

\begin{equation}
Y(s)(1 + KG(s)) = KG(s)R(s)
\end{equation}

\begin{equation}
\frac{Y(s)}{R(s)} = \frac{KG(s)}{1 + KG(s)}
\end{equation}

\vspace{0.3cm}

\textbf{This is the closed-loop transfer function.}

\end{frame}

%------------------------------------------------

\begin{frame}{Important Point}

Poles come from denominator:

\begin{equation}
1 + KG(s) = 0
\end{equation}

This equation depends on $K$

\vspace{0.3cm}

\textbf{Meaning:}

\begin{itemize}
    \item Changing $K$ changes poles
    \item Poles move $\Rightarrow$ root locus
\end{itemize}

\end{frame}

%------------------------------------------------

\begin{frame}{Why output looks like $e^{st}$?}

System is based on differential equations:

\begin{equation}
\frac{dy}{dt} = ay
\end{equation}

Solution:

\begin{equation}
y(t) = e^{at}
\end{equation}

For higher order:

\begin{equation}
y(t) = C_1 e^{s_1 t} + C_2 e^{s_2 t} + \cdots
\end{equation}

\textbf{So poles = values of $s$}

\end{frame}

%------------------------------------------------

\begin{frame}{Stability Meaning}

\begin{equation}
s = -a \Rightarrow e^{-at} \rightarrow 0 \quad (\text{stable})
\end{equation}

\begin{equation}
s = +a \Rightarrow e^{at} \rightarrow \infty \quad (\text{unstable})
\end{equation}

\begin{equation}
s = j\omega \Rightarrow \cos(\omega t) \quad (\text{oscillation})
\end{equation}

\vspace{0.3cm}

\begin{itemize}
    \item Left side $\Rightarrow$ stable
    \item Right side $\Rightarrow$ unstable
    \item Imaginary axis $\Rightarrow$ boundary
\end{itemize}

\end{frame}

%------------------------------------------------

\begin{frame}{Part (a)}

Given poles:

\begin{equation}
(s+1)(s+2)(s+3)
\end{equation}

\begin{equation}
G(s) = \frac{K}{(s+1)(s+2)(s+3)}
\end{equation}

Unity DC gain:

\begin{equation}
G(0) = 1
\end{equation}

\begin{equation}
\frac{K}{1 \cdot 2 \cdot 3} = 1
\end{equation}

\begin{equation}
K = 6
\end{equation}

\begin{equation}
\boxed{G(s) = \frac{6}{(s+1)(s+2)(s+3)}}
\end{equation}

\end{frame}

%------------------------------------------------

\begin{frame}{Part(b)-Root Locus Idea}

Root locus = path traced by CLOSED-LOOP POLES as 
K changes

\begin{equation}
K = 0 \Rightarrow s = -1,-2,-3
\end{equation}

\begin{equation}
K \rightarrow \infty \Rightarrow \text{poles move to infinity}
\end{equation}

\textbf{So poles move as $K$ changes}

\end{frame}

%------------------------------------------------

\begin{frame}{Characteristic Equation}

\begin{equation}
1 + KG(s) = 0
\end{equation}

\begin{equation}
1 + \frac{6K}{(s+1)(s+2)(s+3)} = 0
\end{equation}

\begin{equation}
(s+1)(s+2)(s+3) + 6K = 0
\end{equation}

\begin{equation}
s^3 + 6s^2 + 11s + 6 + 6K = 0
\end{equation}

\end{frame}

%------------------------------------------------

\begin{frame}{Boundary Condition}

Given
for some value of 
K, the closed-loop poles are located at:

\begin{equation}
s = \pm j\sqrt{11}
\end{equation}

\begin{itemize}
    \item Poles on imaginary axis
    \item System oscillates
    \item Last stable point
\end{itemize}

\end{frame}

%------------------------------------------------

\begin{frame}{Finding K}

\begin{equation}
K = \frac{|(s+1)(s+2)(s+3)|}{6}
\end{equation}

\begin{equation}
|j\sqrt{11} + 1| = \sqrt{12}
\end{equation}

\begin{equation}
|j\sqrt{11} + 2| = \sqrt{15}
\end{equation}

\begin{equation}
|j\sqrt{11} + 3| = \sqrt{20}
\end{equation}

\begin{equation}
K = \frac{\sqrt{12 \cdot 15 \cdot 20}}{6}
\end{equation}

\begin{equation}
K = \frac{60}{6} = 10
\end{equation}

\end{frame}

%------------------------------------------------

\begin{frame}{Final Answer}

\begin{equation}
\boxed{K = 10}
\end{equation}

\vspace{0.4cm}

% \textbf{Final understanding:}

% \begin{itemize}
%     \item Poles move when $K$ changes
%     \item Imaginary axis = last safe point
%     \item That $K$ = maximum stable gain
% \end{itemize}

\end{frame}

%------------------------------------------------

\end{document}